\heading{数据结构与算法A-17秋-期中}

\noteinfo{本卷由原试卷 OCR 精修；参考答案按题面逐题推导整理。}

\textbf{一、 选择与填空（每空2 分，共24 分）}
\begin{question}
下面函数的时间复杂度是（    ）\\
\begin{pycode}
def recursive(n, m, k):
    if n <= 0:
        print(m, k)
    else:
        recursive(n - 1, m + 1, k)
        recursive(n - 1, m, k + 1)
\end{pycode}
\begin{choiceoptions}[2]
  \item $O(n*m*k)$
  \item $O(n^2*m^2)$
  \item $O(2^n)$
  \item $O(n!)$
\end{choiceoptions}
\begin{solution}
选 $O(2^n)$。该递归每次把规模为 $n$ 的问题分成两个规模为 $n-1$ 的子问题，除递归调用外只做常数或低阶工作，因此可写成 $T(n)=2T(n-1)+O(1)$。展开递归树，第 $i$ 层有 $2^i$ 个结点，直到深度 $n$，结点总数为 $1+2+\cdots+2^n=O(2^n)$。所以指数项主导，不能化简成多项式复杂度。
\end{solution}
\end{question}

\begin{question}
完成在双循环链表结点p 之后插入s 的操作为(       ):\\
\begin{choiceoptions}[2]
  \item p.next.prev=s; s.prev=p; s.next=p.next; p.next=s;
  \item p.next.prev=s; p.next=s; s.prev=p; s.next=p.next;
  \item s.next=p.next; s.prev=p; p.next.prev=s; p.next=s;
  \item s.prev=p; s.next=p.next; s.prev.next=s; s.next.prev=s;
\end{choiceoptions}
\begin{solution}
A、C、D 的赋值次序均能正确完成在双循环链表结点 $p$ 后插入 $s$。正确插入需要同时维护四条链接：$s$ 的前驱为 $p$，$s$ 的后继为原来的 $p.next$，原后继的前驱改为 $s$，$p$ 的后继改为 $s$。B 错在过早令 \texttt{\detokenize{p.next=s}}，随后 \texttt{\detokenize{s.next=p.next}} 会读到已经被改写的新值，使 \texttt{\detokenize{s.next}} 指向自身，原来的后继结点丢失。
\end{solution}
\end{question}

\begin{question}
设栈S 和队列Q 初始状态为空，元素e1，e2，e3，e4，e5，e6 依次通过栈S，一个元素出栈后即进队列Q，若6 个元素出队序列是e2，e4，e3，e6，e5，e1，则栈S 的容量至少是(       )\\
\begin{choiceoptions}
  \item 2;
  \item 3;
  \item 4;
  \item 6
\end{choiceoptions}
\begin{solution}
容量至少为 3。过程为：压入 $e_1,e_2$ 后弹出 $e_2$；压入 $e_3,e_4$ 后弹出 $e_4,e_3$；压入 $e_5,e_6$ 后弹出 $e_6,e_5,e_1$，最大栈深为 3。
\end{solution}
\end{question}

\begin{question}
设循环队列的容量为40 （序号从0 到39） ，现经过一系列的入队和出队运算后： 1） front=12，rear=19；2）front=19，rear=12；在这两种情况下，循环队列中的元素个数分别为 \blank 和 \blank 。\\
\begin{solution}
按通常约定 \texttt{\detokenize{front}} 指向队首、\texttt{\detokenize{rear}} 指向下一个可插入位置，循环队列元素个数为
\[
(rear-front+40)\bmod40.
\]
把题给两组 \texttt{\detokenize{front}}、\texttt{\detokenize{rear}} 代入即可分别得到 $7$ 和 $33$。公式中加上容量 $40$ 是为了避免 \texttt{\detokenize{rear<front}} 时出现负数，最后再取模回到合法范围。
\end{solution}
\end{question}

\begin{question}
一个n 位的字符串共有 \blank 个子串。\\
\begin{solution}
长度为 $n$ 的字符串中，长度为 $1$ 的子串有 $n$ 个，长度为 $2$ 的有 $n-1$ 个，依此类推，长度为 $n$ 的有 $1$ 个。因此非空子串总数为 $n+(n-1)+\cdots+1=n(n+1)/2$；若教材把空串也计作子串，则再加 $1$。
\end{solution}
\end{question}

\begin{question}
已知二叉树有n 个结点 (n > 0)，则度为2 的结点最多有 \blank 个。\\
\begin{solution}
度为 2 的结点最多为 $\lfloor(n-1)/2\rfloor$。设度为 $0,1,2$ 的结点数分别为 $n_0,n_1,n_2$，二叉树恒有 $n_0=n_2+1$。总结点数 $n=n_0+n_1+n_2=2n_2+n_1+1$，因此 $n_2=(n-n_1-1)/2$。要让 $n_2$ 最大，应使度为 1 的结点数 $n_1$ 尽量小：当 $n$ 为奇数取 $n_1=0$，当 $n$ 为偶数取 $n_1=1$，统一得到 $\lfloor(n-1)/2\rfloor$。
\end{solution}
\end{question}

\begin{question}
用数组存储一个有50 个元素的最小值堆。已知堆中没有重复元素。给出最大元素的下标可能的取值范围为 \blank 到 \blank 。\\
\begin{solution}
若堆数组采用 0 下标，最大元素只能出现在叶结点位置，范围为 $25$ 到 $49$；若采用 1 下标，则为 $26$ 到 $50$。因为这是小根堆，父结点关键字不大于孩子，所以最大关键字不可能出现在任何内部结点，否则它的孩子会更大或等大而违反“最大”的位置判断。含 50 个元素的完全二叉树中，0 下标最后一个内部结点为 $\lfloor(50-2)/2\rfloor=24$，叶子从 $25$ 开始；1 下标最后一个内部结点为 $\lfloor50/2\rfloor=25$，叶子从 $26$ 开始。
\end{solution}
\end{question}

\begin{question}
假设用于通信的电文由5 个字符组成。已知其中一个字符的Huffman 编码为1101，则该Huffman 编码树的平均编码长度为 \blank 。\\
\begin{solution}
仅给出一个 Huffman 码字 \texttt{\detokenize{1101}} 不能唯一确定按频率加权的平均编码长度。平均编码长度需要知道每个字符的频率以及所有字符的码长；一个码字只能说明其中某个叶子的深度为 $4$，不能推出整棵 Huffman 树。若补充约定为 5 个字符等概率，且码长集合为 $1,2,3,4,4$，则平均编码长度为
\[
\frac{1+2+3+4+4}{5}=\frac{14}{5}.
\]
若权值不同，即使码长集合相同，加权平均值也会改变。
\end{solution}
\end{question}

\begin{question}
假设先根次序遍历某棵树的结点次序为 \texttt{GABCDIJEFH}，\\
后根次序遍历该树的结点次序为 \texttt{BADIJCFEHG}。那么 I 结点的父结点是 \blank 。\\
\begin{solution}
$I$ 的父结点是 $C$。先根序第一个结点 $G$ 为根，后根序最后一个结点也为 $G$；递归分割可知 $G$ 的孩子依次为 $A,C,E,H$。在 $C$ 对应的子序列中，先根片段为 $CDIJ$，后根片段为 $DIJC$，所以 $C$ 的孩子为 $D,I,J$，从而 $I$ 的父结点为 $C$。
\end{solution}
\end{question}

\begin{question}
一个拥有144 个结点的完全三叉树，现在从倒数第二层上任取一个结点，该结点为叶结点的概率为 \blank 。\\
\begin{solution}
补全题意为“从倒数第二层任取一个结点，该结点为叶结点的概率”。完全三叉树前 5 层共有 $1+3+9+27+81=121$ 个结点，144 个结点说明第 6 层有 23 个结点，占据第 5 层前 $\lceil23/3\rceil=8$ 个父结点，因此第 5 层叶结点数为 $81-8=73$，概率为 $73/81$。
\end{solution}
\end{question}

\textbf{二、 辨析与简答(共24 分)}
\begin{question}
（6 分）现有一个由单链表和循环单链表构成的特殊链表，单链表的头元素是head，末尾元素为tail，head.….tail 即是一条普通的链表，同时tail 元素还是另一条循环单链表的头元素。如A.B.C.D.E.F.G.H.E 就是一个特殊链表，其中head 为A，tail 为E。\\
现在你只知道head，但是你想知道这个循环链表中有多少元素，但同时你的剩余空间十分的小，所以你想用尽可能小的空间来解决问题，请问你会怎么做？（依据占用空间给分）
\begin{solution}
用 Floyd 快慢指针。先令慢指针每次走一步、快指针每次走两步，若链表有环，两者一定会在环内相遇；若快指针遇到空指针，则不是循环链表。相遇后固定其中一个指针，从相遇点沿 \texttt{next} 继续走，直到再次回到相遇点，走过的边数就是环长；对于纯循环链表，这个环长也就是元素个数。算法只使用常数个指针和计数器，空间复杂度 $O(1)$，时间复杂度 $O(n)$。
\end{solution}
\end{question}

\begin{question}
（6 分）假设以I 和O 分别表示入栈和出栈操作，栈的初始和终态均为空，入栈和出栈的操作序列可表示为仅由I 和O 组成的序列，称可以操作的序列为合法序列，否则为非法序列。如 IOIIOIOO 合法，IOOIOIIO 和IIIOIOIO 为非法。请给出一个算法，判断所给操作序列是否合法。
\begin{solution}
扫描操作串，用计数器记录当前栈深：遇到 I 加 1，遇到 O 减 1。任意前缀中 O 的数量不得超过 I，即栈深不能为负，否则说明在空栈上执行了出栈操作；扫描结束时 I 与 O 数量相等，即栈深为 0，否则说明仍有元素未弹出或出栈次数过多。两条件同时满足则该 I/O 序列合法。若题目还给定栈容量 $m$，则扫描过程中还需额外检查栈深是否超过 $m$。
\end{solution}
\end{question}

\begin{question}
（6 分）给定一棵用数组存储的完全二叉树\{10, 5, 12, 3, 2, 1, 8\}\\
（1） 用筛选建堆法将该完全二叉树调整为最大值堆。画出调整后得到的最大值堆，并说明在调整过程中都依次交换了哪些元素。（注：画堆只需要写出层次遍历序列结果即可）\\
（2） 将6、7、14 按顺序插入（1）得到的最大值堆，堆的空间足够大，画出插入14 后得到的最大值堆，并说明在插入14 的过程中都依次交换了哪些元素。\\
（3） 对（2）得到的堆进行3 次deleteMax，画出第3 次deleteMax 后得到的堆，并说明在第3 次deleteMax 的过程中都依次交换了哪些元素。
\begin{solution}
初始数组 $\{10,5,12,3,2,1,8\}$ 经筛选建最大堆得 $\{12,5,10,3,2,1,8\}$，主要交换 $10\leftrightarrow12$。依次插入 $6,7,14$ 后，插入 14 的交换为 $14\leftrightarrow2$、$14\leftrightarrow7$、$14\leftrightarrow12$，得到 $\{14,12,10,6,7,1,8,3,5,2\}$。随后连续 3 次 \texttt{\detokenize{deleteMax}} 后堆为 $\{8,7,5,6,2,1,3\}$；第 3 次删除中的下滤交换为 $3\leftrightarrow8$、$3\leftrightarrow5$。
\end{solution}
\end{question}

\begin{question}
（6 分）对下列15 个等价对进行合并，给出所得等价类树的图示。在初始情况下，集合中的每个元素分别在独立的等价类中。使用重量权衡合并规则，合并时子树结点少的并入结点数多的那棵（多的那个作为新树根，少的那个根作为新根的直接子结点）；若两棵树规模同样大，则把根值较大的并入根值较小（新树根取值小的）。同时，采用路径压缩优化。\\
(0,2) (1,2) (3,4) (3,1) (3,5) (9,11) (12,14) (12,9) (4,14) (6,7) (8,10) (8,7) (7,11) (10,15) (10,13)
\begin{solution}
按题给重量权衡规则并在查找时路径压缩，所有 $0\sim15$ 最终合并为一个等价类，代表元可取 $0$。重量权衡保证小树挂到大树下，避免树高快速增长；路径压缩则在 \texttt{find} 返回时把访问路径上的结点直接改挂到根。若题目要求画“全部操作结束并完全压缩”的等价类树，可画成以 $0$ 为根，其余结点均为 $0$ 的孩子；若只画严格按操作序列产生的中间形态，应保留那些没有被最后一次 \texttt{find} 访问到的父子边。
\end{solution}
\end{question}

\textbf{三、 算法填空(每空3 分，共24 分)}
\begin{question}
下面的算法利用一个栈将一给定的序列从小到大排序。长度为n的序列由1, 2, …, n这n个连续的正整数组成。请补全下面的代码段，使其可以判断给定的序列是否可以利用一个栈进行排序，如果可以，输出相应的操作。\\
例如：序列4 3 1 2，此序列可以排序，输出：push push push pop push pop pop pop。\\
\begin{pycode}
class Stack:
    def clear(self): ...
    def push(self, item): ...
    def pop(self): ...
    def top(self): ...
    def is_empty(self): ...
    def is_full(self): ...


s = Stack()
sequence = [...]
n = len(sequence)


def is_legal():
    # 判断序列是否可以排序
    for i in range(n):
        for j in range(i + 1, n):
            for k in range(j + 1, n):
                if ________:       # 填空 1
                    return False
    return True


def transform():
    # 输出转换操作
    cur, i = 1, 0
    while ________:                # 填空 2
        while s.is_empty() or ________:    # 填空 3
            ________               # 填空 4
            print("push", end=" ")
            if cur == s.top():
                break
        s.pop()
        print("pop", end=" ")
        cur += 1
\end{pycode}
\begin{solution}
一个排列可由一个栈排成升序，当且仅当不存在 $i<j<k$ 且 $sequence[k]<sequence[i]<sequence[j]$ 的 231 模式。填空为：
  \begin{enumerate}[leftmargin=2em]
    \item \texttt{\detokenize{sequence[k] < sequence[i] < sequence[j]}}；
    \item \texttt{\detokenize{cur <= n}}；
    \item \texttt{\detokenize{cur < s.top()}}；
    \item \texttt{\detokenize{s.push(sequence[i]); i += 1}}。
  \end{enumerate}
\end{solution}
\end{question}

\begin{question}
下面的算法将一个用带右链的先根次序法表示的森林转换为用带度数的后根次序法表示。请利用题目给出的树结点ADT和栈ADT，填充空格，使其成为完整的算法。\\
\begin{pycode}
class RlinkTreeNode:
    def __init__(self, info, ltag, rlink=None):
        self.info = info
        self.ltag = ltag       # 1 表示没有左子结点
        self.rlink = rlink     # 右兄弟


class PostTreeNode:
    def __init__(self):
        self.info = None
        self.degree = 0


post_tree = [PostTreeNode() for _ in range(n)]
count = 0


def convert(node):
    """返回右兄弟链上的结点数量（包括自己）。"""
    global count
    if node.ltag == 1:
        # 填空 5
        post_tree[count].degree = 0
        count += 1
    else:
        tmp = convert(next_node(node))   # 访问第一个子结点
        post_tree[count].info = node.info
        # 填空 6
        count += 1

    if ________:                         # 填空 7
        return convert(node.rlink) + 1
    else:
        # 填空 8
\end{pycode}
\begin{solution}
带右链先根序转带度数后根序的填空：
  \begin{enumerate}[leftmargin=2em]
    \item \texttt{\detokenize{post_tree[count].info = node.info}}；
    \item \texttt{\detokenize{post_tree[count].degree = tmp}}；
    \item \texttt{\detokenize{node.rlink is not None}}；
    \item \texttt{\detokenize{return 1}}。
  \end{enumerate}
算法思想是先递归处理当前结点的所有孩子，再把当前结点写入后根序数组。变量 \texttt{tmp} 统计当前结点实际拥有的孩子个数，因此写入结点信息后还要写入度数。右链 \texttt{rlink} 表示兄弟关系，若兄弟存在，就继续沿右链处理并把兄弟也计入父结点的孩子序列；返回值用于告诉上一层当前子树贡献了一个孩子。这样输出顺序满足“先孩子、后根”，并且每个结点附带正确度数。
\end{solution}
\end{question}

\textbf{四、 算法设计与实现 (14 分)}
\begin{question}
（8 分） 编写算法伪码，对给定的字符串str，返回其最长重复子串及其下标位置。如 str="abcdaccdac"，则子串"cdac"是str 的最长重复子串，下标为2。\\
（重复子串允许重叠）
\begin{solution}
可用后缀数组。构造所有后缀及其起始下标，按字典序排序；最长重复子串一定是相邻后缀的最长公共前缀之一。依次计算相邻后缀 LCP，取最长者及较小起始下标即可。朴素实现时间 $O(n^2\log n)$，若用线性 LCP 算法可降为 $O(n\log n)$ 或更优。
\end{solution}
\end{question}

\begin{question}
（6 分）给出一个算法，求一棵树的树高。可以写出伪代码，需要相应的注释，或者写出详细算法思路。对时间复杂度无具体要求。
\begin{solution}
树高递归算法：若结点为空，返回 $-1$；若为叶子，返回 $0$；否则递归求所有孩子子树的高度，取最大值后加 $1$ 作为当前结点高度。这个定义把高度理解为从该结点到最深叶子的边数，因此叶子高度为 $0$。若教材按“层数”定义树高，则空树为 $0$、叶结点为 $1$，相应地把上述返回值整体加 $1$ 即可。遍历每个结点一次，时间复杂度 $O(n)$。
\end{solution}
\end{question}

\textbf{五、 分析证明题 (共 14 分)}
\begin{question}
（8 分）二叉树的内部路径长度是指所有节点的深度的和。假设将N 个互不相同的随机元素插入一棵空的二叉搜索树。请证明得到的二叉搜索树的内部路径长度期望为 O(NlogN)。
\begin{solution}
随机插入得到的 BST 与快速排序递归树同分布。设 $E[I_N]$ 为内部路径长度期望，根的左子树规模 $i$ 在 $0\sim N-1$ 上等可能，因此
  \[
  E[I_N]=N-1+\frac1N\sum_{i=0}^{N-1}(E[I_i]+E[I_{N-1-i}]),
  \]
  解得 $E[I_N]=O(N\log N)$。
\end{solution}
\end{question}

\begin{question}
（6 分）证明按先根次序周游树获得的序列与其对应二叉树的前序周游获得的序列相同。
\begin{solution}
树的孩子兄弟表示中，树的“第一个孩子”成为对应二叉树的左孩子，“右兄弟”成为右孩子。树的先根遍历先访问根，再依次访问各子树；对应二叉树的前序遍历也是先访问根，再沿左孩子进入第一个子树，并沿右孩子依次访问兄弟子树。因此两者访问序列相同。
\end{solution}
\end{question}
